\solchapter{plasma}

\sol{plasma}{1} With $\lambda_D = \sqrt{\epsilon_0\kB T/ne^2}$: (a) ionosphere $2.2\un{mm}$, $n\lambda_D^3\approx10^4$; (b) fusion plasma $70\un{\mu m}$, $n\lambda_D^3\approx3\times10^7$; (c) solar core $2.7\times10^{-11}\un{m}$, $n\lambda_D^3\approx2$. The first two are good plasmas, dominated by collective effects. The solar core is so dense that each particle has only a couple of neighbours within a Debye length --- a mildly ``strongly coupled'' plasma, where Debye screening is only a rough description.

\sol{plasma}{2} $f_p\approx9\sqrt{10^{12}}\un{Hz} = 9\un{MHz}$. AM waves at $1\un{MHz}$ are below the plasma frequency and are reflected, bouncing between ionosphere and ground far beyond the horizon. (By day, a lower layer absorbs them; it disappears at night.) FM at $100\un{MHz}$ is above $f_p$ and passes into space, so it reaches only as far as line of sight.

\sol{plasma}{3} $r_L = mv/eB = (1.67\times10^{-27})(4\times10^5)/[(1.6\times10^{-19})(5\times10^{-9})]\approx840\un{km}$, about an eighth of the Earth's radius. An electron at the same speed has a radius $1836$ times smaller, about $0.5\un{km}$.

\sol{plasma}{4} For $v\ll c$, a frame moving with velocity $\vb u$ sees the electric field $\vb E' = \vb E + \vb u\times\vb B$. Choose $\vb u = \vb E\times\vb B/B^2$. For $\vb E\perp\vb B$, $(\vb E\times\vb B)\times\vb B = -B^2\vb E$, so $\vb E' = 0$. In that frame there is only a magnetic field, and particles simply gyrate about fixed centres; back in the original frame, their centres drift at $\vb u$. The drift comes from a change of frame, which knows nothing about the particle, so it is the same for every charge and mass. Physically: a positive charge speeds up on the half of its orbit where it moves along $\vb E$ and so has a larger radius there; a negative charge gyrates the other way but is also pushed the other way, and drifts in the same direction.

\sol{plasma}{5} $\omega_p = \sqrt{ne^2/\epsilon_0m_e} = 1.6\times10^{16}\un{rad/s}$, so $\hbar\omega_p\approx11\un{eV}$, in the far ultraviolet. Visible photons ($1.7$--$3.1\un{eV}$) are below it and reflected; X-rays are far above it, where the electrons cannot follow the field, and pass through. In gold, transitions of bound $5d$ electrons into the conduction band begin at about $2.4\un{eV}$, absorbing blue and violet light, so the reflected light is yellow. (Copper's similar threshold, near $2.1\un{eV}$, makes it reddish.)

\sol{plasma}{6} Energy conservation keeps $v$ fixed, and $\mu = mv_\perp^2/2B$ is conserved, so $v_\perp^2 = v^2\sin^2\theta\,B/B_{\min}$. The particle reflects where $v_\perp$ reaches $v$; this happens before $B_{\max}$ if $\sin^2\theta\,B_{\max}/B_{\min}>1$, i.e.\ $\sin^2\theta>B_{\min}/B_{\max}$. For a mirror ratio of $4$, the loss cone is $\theta<30^\circ$ (or $>150^\circ$); for an isotropic distribution the fraction of directions inside it is $1-\cos30^\circ = 13\%$.

\sol{plasma}{7} $v_A = B/\sqrt{\mu_0\rho} = 10^{-3}/\sqrt{(4\pi\times10^{-7})(10^{15})(1.67\times10^{-27})}\approx7\times10^5\un{m/s}$. Crossing $10^8\un{m}$ takes about $150\un{s}$.

\sol{plasma}{8} $\tau_E\ge3\times10^{21}/(10^{20}\times15) = 2\un{s}$. Heat leaks out by diffusion across the field, which takes a time $\sim L^2/D$ for a plasma of size $L$, so confinement improves rapidly with size; also, fusion power grows with volume while losses grow with surface area. This is why ITER is enormous.

\sol{plasma}{9} $v_{\text{phase}} = \omega/k = \sqrt{c^2 + \omega_p^2/k^2}>c$, and $v_{\text{group}} = \dd\omega/\dd k = c^2k/\omega<c$, so their product is $c^2$. The phase velocity exceeds $c$, but it is the speed of a pattern of crests that carries no energy or information; signals travel at the group velocity, which is below $c$. With $E = \hbar\omega$ and $p = \hbar k$, the dispersion relation becomes $E^2 = (\hbar\omega_p)^2 + (pc)^2$: a particle of mass $\hbar\omega_p/c^2$.

\sol{plasma}{10} By Amp\`ere's law the toroidal field falls as $1/R$, so $\grad B$ points inward. The grad-$B$ drift, $\propto\vb B\times\grad B/q$, is vertical and has opposite signs for ions and electrons, so charge separates, say ions up and electrons down. The resulting vertical $\vb E$, crossed with the toroidal $\vb B$, gives an $\vb E\times\vb B$ drift that points outward for every particle, carrying the plasma into the wall within microseconds. Twisting the field lines helically makes each line wind poloidally, visiting the top and bottom and the inside and outside of the torus in turn; electrons flowing along the lines short out the vertical charge separation, and each particle's drifts average away. A sphere won't work because, by the hairy ball theorem, any tangent field on a sphere has a zero somewhere: confining field lines must lie on surfaces with the topology of a torus.

\sol{plasma}{11} In the frame moving with the wave, particles moving slightly slower than $v_\phi$ drift backward slowly and are reflected by the wave's potential hills, gaining speed; those slightly faster are reflected the other way and lose speed. Particles therefore gain energy from the wave in proportion to $f(v_\phi - \delta)$ and give energy to it in proportion to $f(v_\phi+\delta)$, so the net transfer to the wave is $\propto f'(v_\phi)$. For a Maxwellian, $f'<0$ at any $v_\phi>0$: the wave damps. A beam adds a bump on the tail of the distribution, creating a range of speeds where $f'>0$. Waves with phase velocities there grow, drawing energy from the beam until they flatten the bump into a plateau.
